\begin{enumerate}

\item Calcular las reacciones del caso real mediante equilibrio y, en su caso, compatibilidad de flechas.

\item Escribir la ley flectora real M(x) en cada intervalo.

\item Introducir una fuerza virtual unidad en el punto y sentido del desplazamiento buscado.

\item Calcular la ley flectora virtual y sumar las integrales de M por M virtual dividido entre EI para todos los tramos.

\item Verificar el resultado mediante la elástica calculada con constantes de integración y condiciones de apoyo.

\item Liberar el rodillo y anular su flecha: X = P(3L-a)/(2a) cuando L es la longitud TOTAL y a la posición del rodillo.

\end{enumerate}

\noindent\textbf{Operaciones y resultados intermedios.}

\begin{enumerate}

\item Integral M por M virtual dividida por EI, tramo 1: 0 a 3000 mm.

\[\fitmath{y_{1}=\frac{4.21875\times10^{12}}{4.6605\times10^{12}}\approx 0.905214\,\mathrm{mm}}\]

\item Integral M por M virtual dividida por EI, tramo 2: 3000 a 4500 mm.

\[\fitmath{y_{2}=\frac{2.8125\times10^{12}}{4.6605\times10^{12}}\approx 0.603476\,\mathrm{mm}}\]

\item Desplazamiento total por suma de tramos.

\[\fitmath{y_{3}=\frac{4.21875\times10^{12}}{4.6605\times10^{12}}+\frac{2.8125\times10^{12}}{4.6605\times10^{12}}\approx 1.50869\,\mathrm{mm}}\]

\end{enumerate}
